\documentclass[11pt,a4paper]{article}
\usepackage{fontspec,amsmath,amssymb,amsthm,microtype}
\setmainfont{lmroman10-regular.otf}[BoldFont=lmroman10-bold.otf,ItalicFont=lmroman10-italic.otf,BoldItalicFont=lmroman10-bolditalic.otf]
\usepackage{xcolor}
\usepackage[most]{tcolorbox}
\usepackage{mathrsfs}
\providecommand{\lt}{<}
\providecommand{\gt}{>}
\setmonofont{DejaVuSansMono.ttf}[BoldFont=DejaVuSansMono-Bold.ttf,ItalicFont=DejaVuSansMono-Oblique.ttf,BoldItalicFont=DejaVuSansMono-BoldOblique.ttf,Scale=0.86]
\newfontfamily\leanfallback{STIXTwoMath-Regular.otf}[Scale=0.86]
\definecolor{leanrule}{gray}{0.62}
\definecolor{leanback}{gray}{0.962}
\definecolor{leantitle}{gray}{0.22}
\definecolor{leanmuted}{gray}{0.38}
\newtcolorbox{leanbox}[1]{enhanced,breakable,colback=leanback,colframe=leanrule,
  boxrule=0.4pt,arc=0.6mm,left=2.4mm,right=2.4mm,top=1.8mm,bottom=1.8mm,
  before skip=0.8em,after skip=1em,fonttitle=\small\bfseries,coltitle=leantitle,
  colbacktitle=leanback,titlerule=0.3pt,title={#1}}
\newlength{\leancharwidth}
\newenvironment{leancode}{\par\smallskip\begingroup\ttfamily\footnotesize\raggedright
  \settowidth{\leancharwidth}{M}\setlength{\parskip}{0pt}\setlength{\parindent}{0pt}
  \frenchspacing}{\par\endgroup\smallskip}
\newcommand{\leanline}[2]{\hangindent=\dimexpr#1\leancharwidth+4\leancharwidth\relax
  \hangafter=1\noindent\hspace*{#1\leancharwidth}#2\par}
\newcommand{\leansource}[1]{\par{\footnotesize\color{leanmuted}#1}\par}
\newcommand{\leanlabel}[1]{\par\smallskip\noindent{\small\bfseries #1}\par}
\usepackage[margin=25mm]{geometry}
\usepackage[unicode,hidelinks]{hyperref}
\hypersetup{pdfauthor={LangevinLean Project}}
\setlength{\emergencystretch}{3em}
\setlength{\parskip}{0.4em}
\setlength{\parindent}{0pt}
\allowdisplaybreaks[2]
\raggedbottom
\title{\texorpdfstring{Closedness of the conditional state-measurability property}{Closedness of the conditional state-measurability property}}
\author{LangevinLean Project}
\date{}
\begin{document}\maketitle
Draft mathematical exposition; these notes may contain mistakes.


Conditional expectation of a measurable function of a current state and an independent future input depends only on the state. Expressed inside the true \(L^2\) space, this property is closed under strong limits and can therefore pass from actual Euler approximants to the diffusion.


\section{\texorpdfstring{Independent-input averaging and strong-limit closure}{Independent-input averaging and strong-limit closure}}


\textbf{Claim.} If \(Z=F(X,W)\) is square integrable, \(X\) is past-measurable and \(W\) is independent of the whole past, then \(\mathbb E[Z\mid\mathcal F]\) is measurable from \(X\) alone. If \(Z_i\to Z\) strongly and all such conditional means are state-measurable, the limit conditional mean is state-measurable too.



\textbf{Proof.} Square integrability implies integrability on the probability space. The independent-input averaging theorem identifies ordinary conditional expectation as an average over the independent input; the actual \(L^2\) projection has the same almost-everywhere representative. For the limit, conditional expectation is a bounded linear projection, hence continuous. The square-integrable random variables measurable from the state form a closed subspace, so the strong limit remains in it.

\begin{leanbox}{Lean: \texttt{condExpL2\_\allowbreak{}independent\_\allowbreak{}input\_\allowbreak{}state\_\allowbreak{}measurable}, \texttt{condExpL2\_\allowbreak{}state\_\allowbreak{}measurable\_\allowbreak{}of\_\allowbreak{}tendsto}}
\leansource{\texttt{condExpL2\_\allowbreak{}independent\_\allowbreak{}input\_\allowbreak{}state\_\allowbreak{}measurable} \quad WassersteinMixing/\allowbreak{}ConditionalMarkovClosure.lean, lines 26–39}
\begin{leancode}
\leanline{0}{theorem\ condExpL2\_independent\_input\_state\_measurable}
\leanline{4}{(X\ :\ Ω\ →\ A)\ (Y\ :\ Ω\ →\ B)\ (hX\ :\ AEStronglyMeasurable[m]\ X\ Q)\ (hY\ :\ Measurable\ Y)}
\leanline{4}{(hi\ :\ Indep\ ((inferInstance\ :\ MeasurableSpace\ B).comap\ Y)\ m\ Q)}
\leanline{4}{(Z\ :\ Lp\ H\ 2\ Q)\ (F\ :\ A\ ×\ B\ →\ H)\ (hF\ :\ StronglyMeasurable\ F)}
\leanline{4}{(hZ\ :\ Z\ =ᵐ[Q]\ fun\ ω\ ↦\ F\ (X\ ω,\ Y\ ω))\ :}
\leanline{4}{AEStronglyMeasurable[mA.comap\ X]\ (condExpL2\ H\ ℝ\ hm\ Z\ :\ Lp\ H\ 2\ Q)\ Q}
\end{leancode}
\leansource{\texttt{condExpL2\_\allowbreak{}state\_\allowbreak{}measurable\_\allowbreak{}of\_\allowbreak{}tendsto} \quad WassersteinMixing/\allowbreak{}ConditionalMarkovClosure.lean, lines 43–54}
\begin{leancode}
\leanline{0}{theorem\ condExpL2\_state\_measurable\_of\_tendsto\ \{I\ :\ Type*\}\ \{l\ :\ Filter\ I\}\ [l.NeBot]}
\leanline{4}{(X\ :\ Ω\ →\ A)\ (hXa\ :\ Measurable\ X)}
\leanline{4}{(Z\ :\ I\ →\ Lp\ H\ 2\ Q)\ (W\ :\ Lp\ H\ 2\ Q)\ (hlim\ :\ Tendsto\ Z\ l\ ({\leanfallback 𝓝}\ W))}
\leanline{4}{(hZ\ :\ ∀\ i,\ AEStronglyMeasurable[mA.comap\ X]\ (condExpL2\ H\ ℝ\ hm\ (Z\ i)\ :\ Lp\ H\ 2\ Q)\ Q)\ :}
\leanline{4}{AEStronglyMeasurable[mA.comap\ X]\ (condExpL2\ H\ ℝ\ hm\ W\ :\ Lp\ H\ 2\ Q)\ Q}
\end{leancode}
\leanlabel{How the Lean proof follows the argument}
The first theorem identifies \texttt{cond\allowbreak{}Exp\allowbreak{}L2} with ordinary conditional expectation and applies the actual averaging theorem. The second uses continuity of that projection and the closed measurable subspace \texttt{is\allowbreak{}Closed\_\allowbreak{}aestrongly\allowbreak{}Measurable}.
\end{leanbox}

\end{document}
